% !TEX TS-program = pdflatex
% Lettre de motivation — DevSecOps / DevOps — Sofrecom Tunisie
\documentclass[a4paper,11pt]{article}
\usepackage[utf8]{inputenc}
\usepackage[T1]{fontenc}
\usepackage[scaled=0.94]{helvet}
\renewcommand{\familydefault}{\sfdefault}
\usepackage[left=2.5cm,right=2.5cm,top=2cm,bottom=2cm]{geometry}
\usepackage{xcolor}
\usepackage[hidelinks]{hyperref}
\definecolor{navy}{RGB}{23,54,93}
\pagestyle{empty}
\setlength{\parindent}{0pt}
\setlength{\parskip}{0pt}
\hypersetup{pdftitle={Lettre de motivation - Baha Eddine Belhaj Mouldi - Sofrecom DevSecOps},pdfauthor={Baha Eddine Belhaj Mouldi}}
\begin{document}

\begin{flushright}
{\bfseries Baha Eddine Belhaj Mouldi}\\
Tunis, Tunisie\\
+216 55 128 982\\
\href{mailto:bahaeddine.belhajmouldi@gmail.com}{bahaeddine.belhajmouldi@gmail.com}\\
\href{https://linkedin.com/in/bahamouldi}{linkedin.com/in/bahamouldi}\\
\href{https://github.com/bahamouldi}{github.com/bahamouldi}
\end{flushright}

\vspace{0.8cm}

\begin{flushleft}
Sofrecom Tunisie\\
Service Recrutement

\vspace{0.8cm}

Tunis, le 28 juillet 2026

\vspace{0.6cm}

{\bfseries Objet : Candidature spontanée --- Ingénieur DevSecOps / DevOps}
\end{flushleft}

\vspace{0.6cm}

Madame, Monsieur,

Sofrecom accompagne les opérateurs télécoms dans des transformations où la cybersécurité et le cloud ne sont plus des sujets annexes mais des conditions de faisabilité. C'est précisément à cette intersection --- entre l'ingénierie de plateforme et la sécurité --- que j'ai construit mon parcours, et c'est ce qui motive ma candidature.

Diplômé de l'École Nationale d'Ingénieurs de Tunis en juin 2026, j'ai consacré mon projet de fin d'études à la conception de BeeWAF, un pare-feu applicatif web intelligent développé chez Beehive Entreprises. Le système combine 10 041 règles de détection et trois modèles de machine learning pour atteindre un score de 98,2/100 sans aucun faux positif. Ce dernier point a été le plus exigeant du projet : réduire les faux positifs à zéro oblige à comprendre finement ce à quoi ressemble le trafic légitime, bien plus que la détection d'attaques elle-même.

Au-delà de la dimension sécurité, ce projet a surtout été un exercice d'industrialisation. BeeWAF est déployé sur Kubernetes à travers un pipeline Jenkins en huit étapes, avec ArgoCD en approche GitOps, et supervisé par des tableaux de bord ELK et Prometheus/Grafana assortis d'un système d'alerting temps réel. J'y ai intégré 27 modules couvrant l'anti-bot, la DLP, l'anti-DDoS et le virtual patching de 37 CVE, en alignement avec sept référentiels de conformité dont OWASP, PCI DSS, NIST et ISO 27001.

Mon expérience du secteur télécom n'est pas théorique. Lors d'un stage en sécurité réseau chez Tunisie Telecom, j'ai développé un projet de détection orienté SOC reposant sur Elastic, Kibana et des règles Sigma, avec simulations d'attaques scriptées et documentation des workflows d'investigation. J'ai par ailleurs mené, en freelance, l'analyse de vulnérabilités et la revue de code sécurité de huit applications web et API, ce qui m'a appris que la valeur d'un rapport de sécurité tient moins à la liste des failles qu'à la capacité à faire comprendre au développeur pourquoi elles existent.

Je viens également du développement --- Python, FastAPI, Java/Spring Boot, Node.js. Cela compte, je crois, pour un profil DevSecOps : je préfère dialoguer avec les équipes de développement et intégrer la sécurité dans leur chaîne de production plutôt que la leur imposer depuis l'extérieur. C'est aussi ce qui me permet d'être opérationnel sur les sujets d'automatisation, au-delà du seul périmètre sécurité.

Rejoindre Sofrecom Tunisie représenterait pour moi l'opportunité de mettre ces compétences au service de projets d'envergure internationale, tout en continuant à progresser au contact d'équipes expérimentées. Je suis basé à Tunis et disponible immédiatement.

Je me tiens à votre disposition pour un entretien et vous prie d'agréer, Madame, Monsieur, l'expression de mes salutations distinguées.

\vspace{0.8cm}

\begin{flushright}
{\bfseries Baha Eddine Belhaj Mouldi}
\end{flushright}

\end{document}
